Por otra parte, para todo $i$, se tiene $N_{\mathfrak p_i}=0$ ya que
$t\notin\mathfrak p_i$~; por tanto $M_{\mathfrak p_i}=P_{\mathfrak p_i}$~; para
demostrar (3.4.1.3), basta hacerlo sustituyendo allí $M$ por $P$~; por
otra parte, si los dos miembros de (3.4.1.3) son iguales, se deduce de la
misma desigualdad para $P$ y de (3.4.1.5) que necesariamente se tiene
$\operatorname{long}(N/tN)=0$, por tanto $N/tN=0$ y finalmente $N=0$, por el lema
de Nakayama, puesto que $N$ es de tipo finito~; ahora bien, $N=0$ significa que
$t$ es $M$-regular. Por tanto, es posible reducirse al caso $M=P$, es decir,
suponer ya que $t$ es $M$-\emph{regular}. Observemos que esto implica que
$\mathfrak m\notin\operatorname{Ass}(M)$, ya que $t$ no puede anular un
elemento $\neq0$ de $M$. Como $A$ es de dimensión $1$, necesariamente se tiene

$\operatorname{Ass}(M)\subset\bigcup_i\{\mathfrak p_i\}$.

Procedamos entonces por inducción sobre
$n=\sum_i\operatorname{long}(M_{\mathfrak p_i})$. Si $n=0$, necesariamente se tiene
$M_{\mathfrak p_i}=0$ para todo $i$, por tanto $M=0$, ya que ninguno de los
$\mathfrak p_i$ pertenece a $\operatorname{Ass}(M)$~; los dos miembros de
(3.4.1.3) son entonces nulos, y la afirmación $\beta$) de (3.4.1.2) es
trivial. Si $n>0$, el razonamiento del comienzo de la demostración de (3.4.1)
permite suponer además que el $A$-módulo $M$ es \emph{fiel}~: esto implica que
$M_{\mathfrak p_i}\neq0$ para todo $i$ (Bourbaki, \emph{Alg. comm.}, cap.~II,
\S~2, n\textsuperscript{o}~2, cor.~2 de la prop.~4), y por consiguiente
$\operatorname{Ass}(M)=\bigcup_i\{\mathfrak p_i\}$.

Supongamos primero que $n=1$~; solo hay entonces un único ideal primo mínimo
$\mathfrak p$ de $A$,
\oldpage[IV-2]{48}
y decir que $M_{\mathfrak p}$ tiene longitud $1$
significa que $M_{\mathfrak p}$ es isomorfo al cuerpo residual
$k=A_{\mathfrak p}/\mathfrak pA_{\mathfrak p}$ como
módulo sobre $A_{\mathfrak p}$. Por tanto, $M_{\mathfrak p}$ es anulado por
$\mathfrak pA_{\mathfrak p}$, de modo que $\mathfrak p$ es el anulador de $M$
(Bourbaki, \emph{Alg. comm.}, cap.~II, \S~2, n\textsuperscript{o}~4,
fórmula~(9)), lo que implica $\mathfrak p=0$ puesto que se supone que $M$ es fiel~;
por tanto, el anillo $A$ es íntegro. Ahora bien, la hipótesis $M\neq0$ implica
$M/tM\neq0$ por el lema de Nakayama, y por consiguiente
$\operatorname{long}(M/tM)\geqslant1$, lo que demuestra (3.4.1.3) en este caso.
Además, si $\operatorname{long}(M/tM)=1$, $M$ es necesariamente \emph{cíclico}
(Bourbaki, \emph{Alg. comm.}, cap.~II, \S~3, n\textsuperscript{o}~2, cor.~2
de la prop.~4), por tanto isomorfo a un cociente $A/\mathfrak b$~; como es
fiel, necesariamente se tiene $\mathfrak b=0$ y $M$ es isomorfo a $A$~;
como $\operatorname{long}(A/tA)=1$, $tA$ es necesariamente igual al ideal
maximal $\mathfrak m$, y como $A$ es un anillo local íntegro noetheriano,
esto prueba que $A$ es un anillo de valoración discreta (Bourbaki,
\emph{Alg. comm.}, cap.~VI, \S~3, n\textsuperscript{o}~6, prop.~9), cuyo uniformizante es $t$.
Recíprocamente, si $A$ es un anillo de valoración discreta, $t$ su
uniformizante, $\operatorname{long}(M_{\mathfrak p})=1$ y si $t$ es
$M$-regular, $M$ es sin torsión, por tanto isomorfo a un submódulo de $A$
($M$ es de tipo finito), y por consiguiente isomorfo al propio $A$, de donde
$\operatorname{long}(M/tM)=\operatorname{long}(A/tA)=1$.

Supongamos ahora que $n\geqslant2$~; existe entonces una sucesión exacta
\[
0\longrightarrow M'\longrightarrow M\longrightarrow M''\longrightarrow0
\]
con $M'\neq0$, $M''\neq0$ y
$\operatorname{Ass}(M)=\operatorname{Ass}(M')\cup\operatorname{Ass}(M'')$~;
en efecto, si $\operatorname{Ass}(M)$ no se reduce a un solo elemento, esto
se deduce de Bourbaki, \emph{Alg. comm.}, cap.~IV, \S~1,
n\textsuperscript{o}~1, prop.~4~; si, por el contrario, $\operatorname{Ass}(M)$ se reduce a
un solo ideal primo, este último es necesariamente el único ideal primo mínimo
$\mathfrak p$ de $A$~; la hipótesis implica entonces que
$\operatorname{long}(M_{\mathfrak p})\geqslant2$ y basta tomar como $M'$ la
imagen inversa de un submódulo de $M_{\mathfrak p}$ distinto de $0$
y de $M_{\mathfrak p}$. Como $t$ es $M$-regular, $t$ no pertenece a ninguno
de los ideales primos de $\operatorname{Ass}(M)$ (Bourbaki,
\emph{Alg. comm.}, cap.~IV, \S~1, n\textsuperscript{o}~1, cor.~2 de la
prop.~2), por tanto, por la misma razón, $t$ es $M'$-\emph{regular} y
$M''$-\emph{regular}. Esta última propiedad implica por (3.4.1.4) que la sucesión
\[
0\longrightarrow M'/tM'\longrightarrow M/tM
\longrightarrow M''/tM''\longrightarrow0
\]
es exacta~; como, por otra parte, lo mismo ocurre con la sucesión
\[
0\longrightarrow M'_{\mathfrak p_i}\longrightarrow M_{\mathfrak p_i}
\longrightarrow M''_{\mathfrak p_i}\longrightarrow0
\]
para todo $i$, se tiene
\[
\begin{aligned}
\operatorname{long}(M/tM)
&=\operatorname{long}(M'/tM')+\operatorname{long}(M''/tM'')\\
\operatorname{long}(M_{\mathfrak p_i})
&=\operatorname{long}(M'_{\mathfrak p_i})
 +\operatorname{long}(M''_{\mathfrak p_i})
\end{aligned}
\]
y la hipótesis de inducción implica, por tanto, la desigualdad (3.4.1.3). Además
los dos miembros solo pueden ser iguales si las desigualdades análogas para
$M'$ y $M''$ son también igualdades. En virtud de la hipótesis de inducción,
esto equivale a la propiedad $\beta$) para los $\mathfrak p_i$ tales que
$M'_{\mathfrak p_i}\neq0$ o $M''_{\mathfrak p_i}\neq0$~; pero estos ideales son
precisamente aquellos para los que $M_{\mathfrak p_i}\neq0$. C.Q.F.D.

\oldpage[IV-2]{49}
